\documentclass{article}
\usepackage[utf8]{inputenc}
\usepackage{verbatim}
\usepackage{siunitx}
\usepackage{listings}
\usepackage{xcolor}

\usepackage[T1]{fontenc}
\usepackage{fontspec}


\definecolor{codegreen}{rgb}{0,0.6,0}
\definecolor{codegray}{rgb}{0.5,0.5,0.5}
\definecolor{codepurple}{rgb}{0.58,0,0.82}
\definecolor{backcolour}{rgb}{0.95,0.95,0.92}

\definecolor{linkblue}{rgb}{0.1,0.3,0.7}

\lstdefinestyle{mystyle}{
    backgroundcolor=\color{backcolour},   
    commentstyle=\color{codegreen},
    keywordstyle=\color{magenta},
    numberstyle=\tiny\color{codegray},
    stringstyle=\color{codepurple},
    basicstyle=\ttfamily\footnotesize,
    breakatwhitespace=false,         
    breaklines=true,                 
    captionpos=b,                    
    keepspaces=true,                 
    numbers=left,                    
    numbersep=5pt,                  
    showspaces=false,                
    showstringspaces=false,
    showtabs=false,                  
    tabsize=2
}

\lstset{style=mystyle}



\usepackage{setspace}
%\usepackage{pgfplots}
\usepackage{tikz}
\usepackage{mathtools}
\usetikzlibrary{arrows.meta}

\onehalfspacing % better to read

\usepackage{multirow} %multiple row for tabular env.
\usepackage{tabularx} %table
\usepackage{makecell}
\usepackage{booktabs}
\usepackage{amsmath,amssymb,amstext} %formulas ect.
\newcommand{\angstrom}{\textup{\AA}} %the unit
\usepackage{indentfirst} %indest first paragraph

\usepackage{csquotes}

\usepackage[hidelinks]{hyperref} %Links to labeled formulas
\hypersetup{
    colorlinks=true,
    linkcolor=linkblue,
    filecolor=magenta,      
    urlcolor=linkblue,
    citecolor=linkblue,
    pdfpagemode=FullScreen,
    }

\usepackage[edges]{forest}
\usepackage{color, colortbl}
\usepackage{geometry} % more space to write on page
\setlength{\topmargin}{-2.5cm}
\setlength{\oddsidemargin}{0cm}
\setlength{\evensidemargin}{0cm}
\setlength{\textwidth}{17cm}
\setlength{\textheight}{23cm}

\setstretch{1.15}
\setlength{\parskip}{6em plus 0.2em}  % Space between paragraphs

\usepackage{graphicx} % Subfigures
\usepackage{subcaption}

%\usepackage{natbib}
\usepackage{graphicx}
\usepackage{tikz}
\usetikzlibrary{arrows}
\usetikzlibrary{calc}
\usetikzlibrary{angles,quotes}
\usepackage{centernot}

%%%%%%%%%%new stuff i added for bachelors
\newcommand{\draft}{\textcolor{red}{Draft:}}
\newcommand{\temp}{\textcolor{red}{Temp:}}
\newcommand{\source}{\textcolor{red}{[source?]}}
\newcommand{\hm}{\textcolor{red}{[hmm?]}}
\newcommand{\form}{\textcolor{red}{[formulation]}}

%\newcommand{\note}[1]{\textcolor{red}{#1}}
\usepackage{gensymb}
\usepackage{tkz-euclide}

\usepackage{ifthen}
\newcommand{\note}[1]{%
    \ifthenelse{\boolean{show_notes}}{%
        \textcolor{red}{#1}%
    }{}%
}

%add numbering thing on equation manually
\newcommand\numberthis{\addtocounter{equation}{1}\tag{\theequation}}

%\setcitestyle{square}

%set darkmode
%\pagecolor[rgb]{0,0,0}
%\color[rgb]{0.7,0.7,0.7}


\graphicspath{ {./} }


\usepackage{parskip}

\begin{document}


\newboolean{show_notes}

\setboolean{show_notes}{false}




\begin{titlepage}
\begin{center}
\vspace*{4cm}
\huge
{\fontfamily{lmtt} Comparing Algorithms Algorithming the Maximum Algorithms
Comparing Problem}
\vspace*{1cm}
\large

by Author

student number: 123

E-Mail: Mail 
\vspace*{1cm}


Bachelor Thesis

Bachelor of Science

Major in Computer Science
\vspace*{1cm}

First Supervisor: 

Second Supervisor: 

University
\vspace*{1cm}

Date of Submission: Submission date
\end{center}
\end{titlepage}


\newpage




%\include{src/list}




\vspace*{2cm}
\begin{center}
\begin{minipage}{0.9\textwidth}
{\large \hspace{2pt}\textbf{Vorwort}\vspace{1pt}}\\

Das Skript ist der Vorlesung von Georg Loho parallel, wobei die Vorlesung auf
	dem Buch G.Fiscer, B.Springborn, Lineare Algebra, Spektrum Berlin,
	Heidelberg 2025. 

Aussagen sind größtenteils wörlich kopiert aus den Vorlesungfolien genauso wie
	diese aus dem Buch kopiert sind, allerdings nicht zwingend stylistisch
	sowie grammatikalisch Korrekt. Beweise, Prosa und Erklärungen sind
	allerdings nicht enthalten.  Anzumerken ist, dass der Latex-Code
	gänzlich originel und ohne KI erstellt wurde.

Das Projekt ist vorzufinden unter
	\url{https://github.com/Feliks24/LinA2_Script}. Der Latex-Code ist zu
	kompelieren mit xelatex.

\end{minipage}
\end{center}
\newpage


{
\hypersetup{linkcolor=black}
\tableofcontents
}
\newpage

\section{Einführung}

\textit{\small Woche 2}

\textbf{Definition:} Eine Menge $R$ zusammen mit zwei Verknüpfungen
\begin{align*}
	+:R\times R\to R, (a,b)\mapsto a+b\\
	\cdot:R\times R \to R, (a,b)\mapsto a\cdot b
\end{align*}
heißt \textbf{Ring}, wenn folgendes gilt:
\begin{itemize}
\item[\textbf{R1}] $R$ zusammen mit der Addition $+$ ist eine abelsche Gruppe
\item[\textbf{R2}] Die Multiplikation $\cdot$ ist assoziativ
\item[\textbf{R3}] Es gelten das Distributivgesetz
\end{itemize}
\vspace{\parskip}

\textbf{Definition:} Ein Ring $R$ ist ein Integritätsring, falls gilt:
\begin{itemize}
\item[$i)$] $R$ ist kommutativ mit Einselement $1\neq0$
\item[$ii)$] $R$ ist nullteilerfrei
\end{itemize}
\vspace{\parskip}

\textbf{Anmerkung:} Ein Nullteiler für $a\neq0$ sei ein $b\neq0$, sodass
$ab=0$. Beispiel:

$\begin{pmatrix}
	1 & 0\\
	0 & 0
\end{pmatrix}
\begin{pmatrix}
	0 & 0\\
	0 & 1
\end{pmatrix}=0$

\textbf{Definion:} Ein Ring $R$ heißt \textbf{nullteilerfrei}, wenn für
alle $a,b\in R$ aus $a\cdot b=0$ stets $a=0$ oder $b=0$ folgt.

\textbf{Definition:} Ist $R$ ein Ring und $R'\subset R$ eine Teilmenge, so
heißt $R'$ \textbf{Unterring}, wenn $R'$ bezüglich der Addition Untergruppe
ist, und bezüglich der Multiplikation für $a,b\in R'$ auch $a\cdot b\in R'$.

Sind $R$ und $S$ Ringe mit Verknüpfungen $+,\cdot$ und $\oplus,\odot$, so heißt
eine Abbildung $\varphi: R\to S$ ein \textbf{Homomorphismus von Ringen}, wenn
für alle $a,b\in R$ gilt:
\begin{align*}
	\varphi(a+b) = \varphi(a)\oplus\varphi(b)\\
	\varphi(a\cdot) = \varphi(a)\odot\varphi(b)
\end{align*}
Zum Beispiel ist $m\mathbb{Z}\subset\mathbb{Z}$ ein Unterring und
$\mathbb{Z}\to\mathbb{Z}/m\mathbb{Z}, a\mapsto a+m\mathbb{Z}$, ein
Homomorphismus.

\textbf{Definition:} Ist $R$ ein Ring mit Einselement $1\neq 0$, so ist seine
\textbf{Charakteristik} erklärt durch
\begin{align*}
	char(R)&:= 0 \qquad \text{falls } n\cdot 1\neq 0 \text{ für alle } n\geq
	1\\
	&:=min\{n\in\mathbb{N}\backslash {0}:n\cdot 1 = 0\} \qquad \text{sonst}
\end{align*}
\textbf{Lemma:} Ist $K$ Körper, so ist char(K) entweder Null oder eine Primzahl

\textbf{Wiederholung} Komplexe Zahlen: 
\begin{itemize}
\item $\mathbb{C}=\{x+iy: x,y\in\mathbb{R}\}$ wobei $i=\sqrt{-1}$
\item $\mathbb{C}$ ist ein Körper mit Charakteristik 0
\item $|x+iy| = \sqrt{x^2 + y^2}$
\end{itemize}\vspace{\parskip}

\textbf{Wiederholung (gekürzt)} Sei $K$ ein Körper, $t$ eine formale Variable
und $a_0,\dots,a_n\in K$. Dann nennt man den Ausdruck 
\begin{equation*}
	f(t)=a_0t^0+a_1t^1+\dots+a_nt^n
\end{equation*}
ein \textbf{Polynom} in $t$. Sei $f(t)\in K[t]:=\{$Polynome in $t$ mit
Koeffizienten as $K\}$

\textbf{Bemerkung:} Ist $K$ ein Körper, so ist die Menge $K[t]$ der Polynome
übre $K$ zusammen mit den oben definierten Verknüpfungen ein kommutativer Ring
ohne Nullteiler. Weiter gilt
\begin{equation*}
	deg(f\cdot g) = deg\ f+ deg\ g \qquad \text{für } f,g\in K[t]
\end{equation*}

Dabei soll formal $n-\infty = -\infty+m= -\infty + (-\infty) = -\infty$ sein

\textbf{Definition: }Sind $a,b\in R$, so heißt $b$ \textbf{Teiler} von $a$
($b|a$), wenn es ein $c\in R$ gibt mit $a=cb$. 

Ein Element $a\in R$ heißt eine \textbf{Einheit}, wenn es invertierbar ist,
d.h., wenne es $\tilde{a}\in R$ gibt mit $a\tilde{a}= 1$.

Sind $a,b\in R$ gegeben, so heißt ein $c\in R$ \textbf{gemeinsamer Teiler} von
$a$ und $b$, wenn $c|a$ und $c|b$. 

Ein gemeinsamer Teiler $d$ von $a$ und $b$ heißt ein \textbf{größter
gemeinsamer Teiler(ggt)}, wenn $c|d$ für jeden gemeinsamen Teiler $c$.

\textbf{Satz:} Sind $f,g\in K[t]$, und ist $g\neq 0$, so gibt es dazu eindeutig
bestimmte Polynome $q,r\in K[t]$ derart, dass $f=qg+r$ und deg $r$ $<$ deg $g$.

\textbf{Satz:} In den Ringen $R=\mathbb{Z}$ und $R=K[t]$ gibt es zu Elementen
$a,b\in R$, die nicht beide 0 sind, stes einen größten gemeinsamen Teiler $d\in
R\backslash\{0\}$.

\textbf{Relation von Bézout:} Sei $R=\mathbb{Z}$ oder $R=K[t]$. Dann gibt es zu
$a,b\in R\backslash \{0\}$ Koeffizienten $x,y\in R$ derart, dass $ggT(a,b) = xa
+yb$.

\textbf{Lemma:} Ist $\lambda\in K$ eine Nullstelle von $f\in K[t]$, so gibt es
ein eindeutig bestimmtes Polynom $g\in K[t]$ mit $f=(t-\lambda)g$ und es gilt
$deg\ g=(deg\ f)-1$.

\textbf{Korollar 1:} Sei $K$ ein beliebiger Körper, $f\in K[t]$ ein Polynom und
$k$ die Anzahl der Nullstellen von $f$. Ist $f$ vom Nullpolynom verschieden, so
gilt $k\leq deg\ f$.

\textbf{Definition:} Ist $f\in K[t]$ vom Nullpolynom verschieden und
$\lambda\in K$, so heißt $\mu(f;\lambda):=max\{r\in\mathbb{N}: f=(t-\lambda)^rg
\text{ mit } g\in K[t]\}$ die Vielfachheit der Nullstelle $\lambda$ von $f$.

\textbf{Fundamentalsatz der Algebra:} Jedes Polynom $f\in \mathbb{C}[t]$ mit
$deg\ f>0$ hat mindestens eine Nullstelle in $\mathbb{C}$.

\textbf{Korollar:} Jedes Polynom $f\in\mathbb{C}[t]$ zerfällt in
Linearfaktoren, d.h. es gibt $a$ und $\lambda_1, \dots, \lambda_n\in\mathbb{C}$
mit $n=deg\ f$, sodass $f=a(t-\lambda_1)\cdot\dots \cdot (t-\lambda_n)$.

\section{Eigenwerte \& Eigenräume}

\textit{\small Woche 3}

\textbf{Definition:} Sei $F$ ein Endomorphismus des $K$-Vektorraums $V$. Ein
$\lambda\in K$ heißt Eigenwert von $F$, wenn es ein $v\in V$ mit $v\neq0$ gibt,
sodass gilt $F(v) = \lambda v$. 

Jedes vom Nullvektor verschiedene $v\in V$ mit $F(v)=\lambda v$ heißt
Eigenvektor von $F$ (zum Eigenwert $\lambda$).

\textbf{Definition:} Ist $F$ ein Endomorphismus von $V$ und $\lambda\in K$, so
nennen wir $Eig(F;\lambda):=\{v\in V:F(v)=\lambda v\}$ den Eigenraum von $F$
bezüglich $\lambda$.

\textbf{Bemerkung:} 
\begin{itemize}
\item[a)] $Eig(F;\lambda)\subset V$ ist ein Unterektorraum
\item[b)] $\lambda$ ist Eigenwert von $F$ $\Leftrightarrow$
	$Eig(F;\lambda\neq\{0\}$
\item[c)] $Eig(F;\lambda)\backslash\{0\}$ ist die Menge der zu $\lambda$
	gehörigen Eigenvektoren von $F$
\item[d)] $Eig(F;\lambda) = Ker(F-\lambda I_V)$
\item[e)] Sind $\lambda_1,\lambda_2\in K$ verschieden, so ist
	$Eig(F;\lambda_1)\cap Eig(F;\lambda_2)=\{0\}$
\end{itemize}
\vspace{\parskip}

\textbf{Lemma:} Gegeben seien $F\in End(V)$ und Eigenvektoren $v_1,\dots, v_m$
zu paarweise verschiedenen Eigenwerten $\lambda_1,\dots, \lambda_m$. Dann sind
$v_1,\dots,v_m$ linear unabhängig. Insbesondere ist $m\leq dim\ V$.

\textbf{Definition:} Sei $F$ ein Endomorphismus und $\mathcal{A}$ eine Basis von
$V$. Ist $A=M_\mathcal{A}(F)$, so nnenen wir $P_F(t):=P_A(t)$ das
charakteristische Polynom vom $F$.

\textbf{Satz:} Sei $V$ ein $K$-Vektorraum der Dimension $n<\infty$ und $F$ ein
Endomorphismus von $V$. Dann hat das charakteristische Polynom $P_F(t)\in K[t]$
die folgenden Eigenschaften:

\begin{enumerate}
	\item $deg\ P_F=n$
	\item $P_F$ beschreibt die charakteristische Funktion $\tilde{P}_F:K\to
		K, \lambda\mapsto det(F-\lambda I_V)$
	\item Die Nullstellen von $P_F$ sind die Eigenwerte von F
	\item Ist $A$ eine Matrix, die $F$ darstellt, so ist $P_F(t)=det(A-t E_n)$
\end{enumerate}
\vspace{\parskip}

\textbf{Satz:} Sei $V$ ein endlichdimensionaler $K$-Vektorraum und $F\in
End(V)$. Dann sind die folgenden Bedingungen äquivalent:

\begin{itemize}
\item[i)] $F$ ist diagonalisierbar
\item[ii)]   	\begin{itemize}
			\item[a)] Das charakteristische Polynom $P_F$ zerfällt
				in Linearfaktoren und
			\item[b)] $dim\ Eig(F;\lambda)=\mu(P_F;\lambda)$ für
				alle Eigenwerte $\lambda$ von $F$, d.h., die
				geometrische Vielfachheit ist gleich der
				algebraischen Vielfachheit
		\end{itemize}
\item[iii)] Sind $\lambda_1,\dots,\lambda_k$ die paarweise verschiedenen
	Eigenwerte von $F$, so ist $V = Eig(F;\lambda_1)\oplus\dots\oplus
	Eig(F;\lambda_k)$
\end{itemize}
\vspace{\parskip}

\textbf{Definition:} Ein Endomorphismus heißt diagonalisierbar, wenn es eine
asis aus Eigenvektoren gibt. $A=T^{-1} D T$

\textbf{Bemerkung:} Ist $dim\ V=n<\infty$, so ist $F\in End(V)$ genau dann
diagonalisierbar, wenn es eine Basis $\mathcal{B}=(v_1,\dots,v_n)$ von $V$
gibt, sodass $M_\mathcal{B}(F)$ eine Diagonalmatrix ist, d.h.
\begin{equation*}
	M_\mathcal{B}(F)=
	\begin{pmatrix}
		\lambda_1 & & 0\\
		& \ddots & \\
		0 & & \lambda_n
	\end{pmatrix}
\end{equation*}

\textbf{Definition:} Sei $F:V\to V$ ein Endomorphismus und $\lambda$ ein
Eigenwert von $F$. Dann heißt $\mu(P_F;\lambda)$ die algebraische Vielfachheit
von $\lambda$ und $dim\ Eig(F;\lambda)$ die geometrische Vielfachheit von
$\lambda$.

\textbf{Lemma:} Ist $\lambda$ Eigenwert von $F$, so gilt $1\leq dim\
Eig(F;\lambda)\leq \mu(P_F;\lambda)$.

\textbf{Definition:} Ein Vektorraum $V$ heißt direkte Summe von
Untervektorräumen $W_1,\dots, W_k$, (in Zeichen $V=W_1\oplus\dots\oplus W_k$),
wenn folgende Bedingungen erfüllt sind:

\begin{itemize} 
\item[DS1] $V=W_1+\dots W_k$
\item[DS2] Sind $w_1\in W_1,\dots,w_k\in W_k$ gegeben mit $w_1+\dots w_k=0$, so
	folgt $w_1=\dots=w_k=0$
\end{itemize}
\vspace{\parskip}

\textbf{Satz:} Für Untervektorräume $W_1,\dots,W_k$ eines endlichdimensionalen
Vektorraums $V$ sind folgende Bedingungen äquivalent:

\begin{itemize}
\item[i)] $V=W_1\oplus\dots\oplus W_k$

\item[ii)] ist für jedes $i\in\{1,\dots,k\}$ eine Basis $(v_1^{(i)},\dots,
	v_{r_i}^{(i)})$ von $W_i$ gegeben, so ist $\mathcal{B}:=(v_{1}^{(1)},
	\dots, v_{r_1}^{(1)},\dots , v_{1}^{(k)}, \dots , v_{r_k}^{(k)})$ eine
	Basis von $V$

\item[iii)] $V=W_1+\dots+W_k$ und $dim\ V=dim\ W_1+ \dots+dim\ W_k$
\end{itemize}
\vspace{\parskip}

\textbf{Satz:} Sei $V$ ein endlichdimensionaler $K$-Vektorraum und $F\in
End(V)$. Dann sind die folgenden Bedingungen äquivalent:

\begin{itemize}
\item[i)] $F$ ist diagonalisierbar
\item[ii)]
	\begin{itemize}
	\item[a)] Das charakteristische Polynom $P_F$ zerfällt in
		Linearfaktoren und
	\item[b)] $dim\ Eig(F;\lambda)=\mu(P_F;\lambda)$ für alle Eigenwerte
		$\lambda$ von $F$
	\end{itemize}
\item[iii)] Sind $\lambda_1,\dots, \lambda_k$ die paarweise verschiedenen
	Eigenwerte von $F$, so ist $V=Eig(F;\lambda_1)\oplus\dots\oplus
	Eig(F;\lambda_k)$
\end{itemize}

\vspace{\parskip}

\section{Trigonalisierbarkeit}

\textit{\small Woche 4}

\textbf{Definion:} Sei $F\to V$ ein Endomorphismus und $W\subset V$ ein
Untervektorraum. $W$ heißt $F$-invariant, wenn $F(W)\subset W$.

\textbf{Bemerkung:} Ist $W\subset V$ ein $F$-invarianter Unterraum, so ist
$P_{F|W}$ ein Teiler von $P_F$.

\textbf{Definition:} Unter einer Fahne $(V_r)$ in einem $n$-dimensionalen
Vektorraum $V$ versteht man eine Kette $\{0\} =V_0\subset
V_1\subset\dots\subset V_n=V$ von Untervektorräumen mit $dim\ V_r=r$. Ist $F\in
End(V)$, so heißt die Fahne $F$-invariant, wenn $F(V_r)\subset V_r$ für alle
$r\in\{0,\dots,n\}$

\textbf{Definition:} Ein Endomorphismus heißt trigonalisierbar, wenn es eine
obere Dreiecksmatrix als Darstellungsmatrix gibt. Eine Matrix heißt
trigonalisierbar, wenn sie ähnlich ist zu einer oberen Dreiecksmatrix.

\textbf{Bemerkung:} Für $F\in End(V)$ sind folgende Bedingungen gleichwertig:
\begin{itemize}
\item[i)] Es gibt eine $F$-invariante Fahne in V
\item[ii)] Es gibt eine Basis $\mathcal{B}$ von $V$, sodass $M_\mathcal{B}(F)$
	obere Dreiecksmatrix ist
\item[iii)] $F$ ist trigonalisierbar
\end{itemize}
\vspace{\parskip}

\textbf{Trigonalisierungssatz:} Für einen Endomorphismus $F$ eines
$n$-dimensionalen $K$-Vektorraums sind folgende Bedingungen äquivalent:
\begin{itemize}
\item[i)] $F$ ist trigonalisierbar
\item[ii)] Das charakteristische Polynom $P_F$ zerfällt in Linearfaktoren, d.h.
	\begin{equation*}
		P_F=\pm(t-\lambda_1)\cdot\dots\cdot(t-\lambda_n)\qquad\text{mit
		}\lambda_1,\dots,\lambda_n\in K
	\end{equation*}
\end{itemize}

\textbf{Korollar:} Jeder Endomorpismus eines endlichdimensionalen komplexen
Vektorraums ist trigonalisierbar.

\section{Jordan'sche Normalform}
\textit{\small Woche 5}

\textbf{Definition:} Ein Jordan-Block der Größe zum Eigenwerte
$\lambda$ ist die quadratische $n\times n$-Matrix

\begin{equation*}
	J_n(\lambda):=
\begin{pmatrix}
	\lambda & 1 & & & & 0\\
	& \lambda & 1 & & &  \\
	& & \ddots & \ddots & &  \\
	& &  & & \lambda & 1 \\
	0 & &  & & & \lambda\\
\end{pmatrix}
\end{equation*}

\textbf{Definition:} Eine Jordan-Matrix ist eine Blockdiagonalmatrix mit
Jordan-Blöcken in der Diagonalen, d.h. eine Matrix der Form

\begin{equation*}
	\begin{pmatrix}
		J_{r_1}(\lambda_1) & & & 0\\
		& J_{r_2}(\lambda_2)  & & \\
		& & \ddots &\\
		0 & & & J_{r_n}(\lambda_n)\\
	\end{pmatrix}
\end{equation*}

wobei die Eigenwerte $\lambda_1,\dots,\lambda_l$ der Jordan-Blöcke nicht
verschieden sein müssen. Man sagt auch, eine Matrix habe Jordansche Normalform,
wenn sie eine Jordan-Matrix ist.

Die Summe der Größen der Jordan-Blöcke zu einem Eigenwerte $\lambda$ ist also
die algebraische Vielfachheit des Eigenwerts:
$\mu(P_F;\lambda)=\sum_{j:\lambda_j=\lambda}r_j.$

Die Anzahl der Jordan-Blöcke zu einem Eigenwert $\lambda$ ist also die
geometrische Vielfachheit des Eigenwerts: $dim\
Eig(F;\lambda)=\sum_{j:\lambda_j=\lambda}1$.

\textbf{Anmerkung:} Somit sind Diagonalmatrizen Spezialfälle von
Jordan-Matrizen, wobei jeder Jordan-Block jeweils die Größe $1\times1$ hat.

\textbf{Satz von der Jordanschen Normalform:} Sei $V$ ein endlichdimensionaler
Vektorraum übre einem algebraisch abgeschlossenen Körper, und sei $F\in
End(V)$. Dann igbt es eine Basis $\mathcal{B}$ von $V$, bezüglich der die
darstellende Matrix $M_\mathcal{B}(F)$ Jordansche Normalform hat. Diese
Jordan-Matrix $M_\mathcal{B}(F)$ ist bis auf die Reihenfolge der Jordan-Blöcke
durch den Endomorphismus $F$ eindeutig bestimmt.

\textbf{Satz von der Jordanschen Normalform (Matrix):} Sei $K$ ein algebraisch
abgeschlossener Körper, und sei $A\in M(n\times x;K)$. Dann ist $A$ ähnlich zu
einer Matrix $J$ in Jordanscher Normalform. Das heißt, es gibt eine Matrix
$T\in GL(n;K)$, sodass $J=TAT^{-1}$ eine Jordan-Matrix ist. Diese Jordan-Matrix
$J$ ist bis auf die Reihenfolge der Jordan-Blöcke durch die Matrix $A$
eindeutig bestimmt.

\textbf{Satz:} Quadratische Matrizen über einem algebraisch abgeschlossenen
Körper sind genau dann ähnlich, wenn sie(bis auf die Reihenfolge der
Jordan-Blöcke) die gleiche Jordansche Normalform haben.

\textbf{Anmerkung:} Ein algebraisch abgeschlossener Köper $K$ ist ein Körper
für welchen gilt: $K[t]$ zerfällt in Linearfaktoren. Zum Beispiel $\mathbb{C}$
aber nicht $\mathbb{R}$

\textbf{Definition:} Gegeben sei $K$-Vektorraum $V$ und $F\in End(V)$ dann sei
\begin{itemize}
\item[i)] ein Hauptvektor erste Stufe zum Eigenwert $\lambda$ der Eigenvektor
	von $F$ zum Eigenwert $\lambda$
\item[ii)] ein Hauptvektor der Stufe $k\geq2$ zum Eigenwert $\lambda$ ein
	Vektor $v\in V$, sodass $F(v)=\lambda v+ v'$ wobei $v'$ ein Hauptvektor
	der Stufe $k-1$ zum Eigenwert $\lambda$ ist
\end{itemize}

\textbf{Lemma:} Ein Vektor $v\in V$ ist genau dann Hauptvektor der Stufe
$k\geq1$ zum Eigenwert $\lambda$, wenn $(F-\lambda I_V)^k(v)=0$ aber
$(F-\lambda I_V)^{k-1}(v)\neq0$.

\textbf{Definition:} Eine Familie $(v_1,\dots, v_k)$ heißt eine Jordan-Kette
von $F$ zum Eigenwert $\lambda$, wenn $v_k$ ein Hauptvektor der Stufe $k$ ist
und für $j\in\{1,\dots,k-1\}$ gilt $v_j=(F-\lambda I_V)(v_{j+1})$. Insbesondere
ist dann $v_j$ ein Hauptvektor der Stufe $j$ zum selben Eigenwert $\lambda$.

\textbf{Satz von der Jordanschen Normalform (Jordan-Ketten):} Sei $V$ ein
endlichdimensionaler Vektorraum über einem algebraisch abgeschlossenen Körper,
und sei $F\in End(V)$. Dann gibt es eine Basis $\mathcal{B}$ von $V$, die aus
Jordan-Ketten von $F$ besteht. Dabei sind die Anzahlen der Jordan-Ketten zu
jedem Eigenwert und ihre Längen durch den Endomorphismus $F$ eindeutig
bestimmt.

\textbf{Algorithmus:} Bestimmung der Jordanschen Normalform
\begin{enumerate}
	\item Eigenwerte bestimmen
	\item Für jeden Eigenwert
		\begin{enumerate}
			\item bestimme linear unabhängige Hauptvektoren der
				höchsten Stufe
			\item berechne zugehörige Jordan-Kette
		\end{enumerate}
\end{enumerate}

\textbf{Anmerkung:} Zu jedem Eigenwert $\lambda$ exestiert ein $l$, sodass alle
Hauptvekotren in $Ker\ (F-\lambda I_V)^l$ wobei sich der Kern für $l+1$ nicht
verändert. Zu bemerken ist, dass $l$ dann die Größe des größten Jordan-Blocks
ist zum Eigenwert $\lambda$.

\section{Polynome und Endomorphismen}

\textit{\small Woche 6}

\textbf{Definition:} Sei $p=a_0+a_1t+\dots a_dt^d\in K[t]$ und $F\in End(V)$
dann gilt $p(F):=a_0I_V+a_1F+\dots+a_dF^d\in End(V)$ wobei $F^k=F\circ\dots
\circ F$. Sei weiterhin $A\in M(n\times n,K)$ dann sei $p(A):= a_0 E_n
+a_1A+\dots+a_dA^d\in M(n\times n,K)$.

\textit{\small Woche 7}

\textit{\small Woche 8}

\textit{\small Woche 9}

\textit{\small Woche 10}

\textit{\small Woche 11}

\textit{\small Woche 12}

\end{document}

